\section{Projeter le pilotage d'équipe dans un contexte professionnel complet}

Le cadre pédagogique de Lynx Immo m'a permis de coordonner une équipe réelle, mais il ne reproduisait pas toutes les responsabilités associées à un pilotage d'équipe en entreprise. Deux limites sont particulièrement importantes au regard du référentiel : l'équipe avait été constituée avant le démarrage du projet et aucun processus RH de recrutement, de gestion de carrière ou d'évaluation individuelle formalisée n'était disponible. Plutôt que de présenter ces dimensions comme réalisées, je peux en déduire le dispositif que je mettrais en place dans un contexte professionnel comparable.

\subsection{Plan de constitution d'une équipe projet}

La première étape consisterait à traduire le périmètre du produit en besoins de compétences, puis à distinguer les compétences indispensables au démarrage de celles pouvant être acquises en cours de projet. Pour Lynx Immo, le noyau minimal aurait dû couvrir le pilotage produit, le développement mobile Flutter, l'administration web React, le backend et les données Supabase/PostgreSQL, l'UX/UI, les tests et la sécurité. Les sujets plus spécialisés --- abonnements mobiles, publication sur les stores ou analyse avancée des données --- auraient pu être attribués à un membre déjà présent, faire l'objet d'une montée en compétences ciblée ou justifier un renfort ponctuel.

Dans une entreprise, je formaliserais cette préparation avec une fiche synthétique par rôle : mission principale, responsabilités, compétences techniques, compétences relationnelles attendues, niveau d'autonomie recherché et interactions avec les autres fonctions. Le plan de constitution pourrait ensuite être partagé avec les RH et le responsable métier afin de choisir entre mobilité interne, formation et recrutement externe. L'objectif ne serait pas de multiplier les profils, mais d'éviter les dépendances critiques et de s'assurer que chaque domaine indispensable possède un responsable identifié et une solution de continuité.

\begin{table}[H]
\centering
\small
\begin{tabular}{p{3.4cm} p{4.4cm} p{3.3cm} p{3.0cm}}
\toprule
\textbf{Domaine} & \textbf{Responsabilité attendue} & \textbf{Mode de couverture privilégié} & \textbf{Point de vigilance} \\
\midrule
Pilotage produit & Périmètre, priorités, relation client, risques et arbitrages & Responsable de projet identifié dès le lancement & Éviter le cumul excessif avec la réalisation technique \\
Mobile & Parcours apprenant, intégration backend, qualité Android/iOS & Compétence interne Flutter ou formation ciblée & Prévoir un relais sur le socle applicatif \\
Web administrateur & Interface de gestion, usages métier, intégration API & Développeur React/TypeScript & Synchroniser les contrats de données avec le backend \\
Backend et données & Modèle, authentification, RLS, fonctions serveur, migrations & Profil backend/data ou référent technique & Revue croisée des règles d'accès \\
UX/UI & Parcours, maquettes, cohérence visuelle et accessibilité & Compétence dédiée ou contribution transverse formalisée & Faire valider tôt les parcours prioritaires \\
Tests et sécurité & Stratégie de validation, non-régression, scénarios sensibles & Responsabilité transverse explicite & Ne pas reporter ces activités à la fin du projet \\
\bottomrule
\end{tabular}
\caption{Exemple de plan de constitution cible pour une équipe de type Lynx Immo}
\label{tab:lynx-plan-constitution-cible}
\end{table}

Ce plan constitue une projection professionnelle et non la description du groupe étudiant réellement constitué. Il montre néanmoins comment les constats du projet --- concentration des responsabilités, disponibilité inégale et dépendance à quelques membres --- peuvent être transformés en critères concrets de constitution d'équipe.

\subsection{Construire un référentiel de performance utilisable pendant le projet}

Le suivi réalisé pendant Lynx Immo reposait surtout sur l'avancement des exigences, Jira, les jalons, les échanges réguliers, la rétrospective finale et les entretiens individuels. Ces éléments donnent une vision utile du fonctionnement collectif, mais ils ne constituent pas à eux seuls un référentiel de performance individuel et collectif au sens d'un dispositif RH structuré.

Dans un contexte professionnel, je définirais avant les premières itérations un petit nombre de critères observables, adaptés aux responsabilités de chacun. Pour l'équipe, les critères pourraient porter sur la tenue des objectifs de sprint, la qualité des livrables, la résolution des blocages, la stabilité du produit et la capacité à partager l'information. Pour chaque membre, le référentiel ne mesurerait pas uniquement le volume de tickets terminés : il intégrerait la qualité, l'autonomie, la collaboration, la fiabilité des engagements et la progression sur les compétences attendues du rôle.

\begin{table}[H]
\centering
\small
\begin{tabular}{p{3.2cm} p{5.2cm} p{5.2cm}}
\toprule
\textbf{Dimension} & \textbf{Collectif} & \textbf{Individuel} \\
\midrule
Livraison & Objectifs de sprint atteints, exigences réellement terminées, défauts bloquants & Engagements tenus et qualité des livrables confiés \\
Qualité & Non-régression, dette identifiée, revues effectuées & Prise en compte des revues, tests et règles communes \\
Collaboration & Blocages partagés rapidement, décisions tracées, entraide & Communication, feedback, capacité à demander ou apporter de l'aide \\
Autonomie & Réduction des dépendances critiques, partage de connaissance & Capacité à avancer dans son domaine et à expliciter ses choix \\
Progression & Compétences manquantes identifiées et actions engagées & Acquisition ou consolidation des compétences prévues pour le rôle \\
\bottomrule
\end{tabular}
\caption{Référentiel de performance cible pour le suivi d'une équipe projet}
\label{tab:lynx-referentiel-performance-cible}
\end{table}

Je compléterais ce référentiel par des feedbacks courts et réguliers, par exemple à la fin de plusieurs sprints, plutôt que d'attendre la rétrospective finale. Chaque échange permettrait de comparer les objectifs du rôle, les résultats observés, les difficultés rencontrées et les besoins de développement. Le feedback serait à double sens : le responsable doit pouvoir signaler un écart ou valoriser une progression, mais le membre de l'équipe doit également pouvoir exprimer une charge excessive, un manque de clarté ou une difficulté liée à l'organisation.

Lorsque l'écart provient d'une compétence insuffisamment maîtrisée, la réponse ne devrait pas être uniquement évaluative. Elle peut conduire à une revue en binôme, une formation ciblée, un prototype limité, une documentation partagée ou une réaffectation temporaire. Dans un cadre RH réel, les résultats de ces échanges pourraient ensuite alimenter des objectifs de développement plus longs et, lorsque cela est pertinent, les perspectives d'évolution du collaborateur.

Cette projection ne remplace pas l'expérience RH que je n'ai pas eue pendant le projet. Elle explicite en revanche ce que les limites observées m'ont appris : constituer une équipe ne consiste pas seulement à répartir des noms sur des tâches, et évaluer sa performance ne consiste pas à compter des tickets. Le pilotage doit relier besoins de compétences, responsabilités, résultats observables, feedback et développement des personnes. Cette démarche complète ainsi les preuves réelles de C4.2 et C4.6 par un plan professionnel cohérent avec les situations rencontrées.